\chapter{Smart contracts et ingénierie crypto}\label{ch:iv3:smart-contracts-and-crypto-engineering}

\begin{ivtrains}
Ce qu'un ingénieur en trading d'une entreprise crypto doit savoir~: comment l'EVM facture le travail effectué, les failles classiques des contrats et leurs correctifs, les oracles, la MEV, et l'infrastructure d'une bourse. Chaque contrat vulnérable est présenté avec sa faille et sa correction, sur un EVM local en mémoire ou un modèle Python, et aucun ne fait référence à un déploiement réel. Ici, Solidity est en version 0.8.28, compilé par solc pour le fork Cancun, non optimisé.
\end{ivtrains}

\begin{question}[$\star$ \iqroles{QD, Dev} \iqfirm{crypto} \iqlevel{junior}]\label{iq:iv3:smart-contracts-and-crypto-engineering:1}
Alice détient deux coins non dépensés, valant 500 et 300 unités, sur une chaîne qui utilise des sorties de transaction non dépensées. Elle paie 600 avec des frais de 10. En Python, écrivez la dépense~: ce qui entre, ce qui sort, et ce qu'elle détient ensuite~? Comment ce même paiement se présenterait-il sur une chaîne à comptes~? Nommez une chose que chaque modèle rend plus facile.
\end{question}

\begin{question}[$\star$ \iqroles{QD, Dev, QT} \iqfirm{crypto} \iqlevel{junior}]\label{iq:iv3:smart-contracts-and-crypto-engineering:2}
Sur Ethereum, un transfert simple consomme 21\,000 gas. Les frais de base sont de 20 gwei, et vous ajoutez des frais de priorité de 1 gwei. Un gwei vaut $10^{-9}$ ETH. Quel est le montant des frais en ETH, et qui reçoit chaque partie~? Si le bloc est plein, de combien les frais de base du bloc suivant peuvent-ils augmenter~?
\end{question}

\begin{question}[$\star$ \iqroles{QD, Dev} \iqfirm{crypto} \iqlevel{junior}]\label{iq:iv3:smart-contracts-and-crypto-engineering:3}
Dans le \cref{lst:iv3:cr:fixed}, en Solidity 0.8, que renvoient \texttt{divFirst(1, 3, 3)}, \texttt{mulFirst(1, 3, 3)}, \texttt{wmul(1.5e18, 2e18)} et \texttt{wmul(1, 1)}~? Que fait \texttt{mulFirst} sur $(2^{255}, 2, 2)$, alors que la vraie réponse tient sur 256 bits~? Pourquoi Solidity n'a-t-il pas de virgule flottante, et comment écrivez-vous 1.5 à la place~?
\omcode{code/interviews-3/08-smart-contracts-and-crypto-engineering/solidity/Fixed.sol}{4}{18}{Arithmétique entière.\protect\hfil\protect\break}\label{lst:iv3:cr:fixed}
\end{question}

\begin{question}[$\star$ \iqroles{QD, Dev} \iqfirm{crypto} \iqlevel{junior}]\label{iq:iv3:smart-contracts-and-crypto-engineering:4}
Un nouveau déploiement du \cref{lst:iv3:cr:locations} a \texttt{p.x} égal à 1. Que renvoie \texttt{viaMemory()}, et que renvoie \texttt{viaStorage()} ensuite~? Que signifient \texttt{memory}, \texttt{storage} et \texttt{calldata}, et combien de temps dure chacun~? Lequel coûte le plus cher~?
\omcode{code/interviews-3/08-smart-contracts-and-crypto-engineering/solidity/Locations.sol}{4}{19}{Emplacements des données.\protect\hfil\protect\break}\label{lst:iv3:cr:locations}
\end{question}

\begin{question}[$\star$ \iqroles{QD, Dev} \iqfirm{crypto} \iqlevel{junior}]\label{iq:iv3:smart-contracts-and-crypto-engineering:5}
Qu'est-ce que le nonce d'un compte sur Ethereum~? Qu'est-ce qui empêche quelqu'un de rediffuser votre transfert signé une seconde fois~? Le prochain nonce de votre compte est 0, et vos transactions de nonces 0, 1 et 3 sont en attente. Lesquelles peuvent être incluses, et que devient la dernière~? Que se passe-t-il si deux transactions partagent un même nonce~?
\end{question}

\begin{question}[$\star$ \iqroles{QD, Dev, QT} \iqfirm{crypto} \iqlevel{junior}]\label{iq:iv3:smart-contracts-and-crypto-engineering:6}
Pourquoi une bourse ne crédite-t-elle pas un dépôt dès qu'elle le voit dans un bloc~? Sur Ethereum en preuve d'enjeu, avec des créneaux de 12 secondes et 32 créneaux par époque, la finalité prend environ deux époques. Combien de temps cela représente-t-il~? Combien coûte à un attaquant la réversion d'un bloc finalisé, et qu'en est-il d'un bloc qui n'a qu'un créneau de profondeur~?
\end{question}

\begin{question}[$\star\star$ \iqroles{QD, Dev} \iqfirm{crypto} \iqlevel{junior, confirmé}]\label{iq:iv3:smart-contracts-and-crypto-engineering:7}
Dans le \cref{lst:iv3:cr:slot}, selon les règles de gas du fork Cancun, le programme envoie quatre transactions distinctes~: \texttt{set(1)}, \texttt{set(2)}, \texttt{set(2)} et \texttt{set(0)}, en partant de \texttt{v} égal à 0. Le surcoût propre à l'appel est le même à chaque fois. En retirant ce surcoût, combien coûte chaque écriture en gas, et y a-t-il un remboursement~? Pourquoi les coûts sont-ils ce qu'ils sont~?
\omcode{code/interviews-3/08-smart-contracts-and-crypto-engineering/solidity/Slot.sol}{4}{10}{Un emplacement de stockage.\protect\hfil\protect\break}\label{lst:iv3:cr:slot}
\end{question}

\begin{question}[$\star\star$ \iqroles{QD, Dev} \iqfirm{crypto} \iqlevel{junior, confirmé}]\label{iq:iv3:smart-contracts-and-crypto-engineering:8}
Combien d'emplacements de stockage chaque contrat présenté dans le \cref{lst:iv3:cr:packing} utilise-t-il~? Dessinez l'emplacement 0 de \texttt{Tight}. Quelle économie cela représente-t-il au déploiement~? Dans quels cas ne réordonneriez-vous pas vos variables~? Pourquoi ne pouvez-vous pas simplement les réordonner dans une nouvelle version d'un contrat évolutif~?
\omcode{code/interviews-3/08-smart-contracts-and-crypto-engineering/solidity/Packing.sol}{4}{14}{Deux agencements.\protect\hfil\protect\break}\label{lst:iv3:cr:packing}
\end{question}

\begin{question}[$\star\star$ \iqroles{QD, Dev} \iqfirm{crypto} \iqlevel{junior, confirmé}]\label{iq:iv3:smart-contracts-and-crypto-engineering:9}
Que font \texttt{checkedDec(0)}, \texttt{uncheckedDec(0)} et \texttt{narrow(300)} dans le \cref{lst:iv3:cr:overflow}~? Lesquelles le compilateur protège-t-il, et lesquelles non~? Pourquoi écrirait-on \texttt{unchecked}~? Donnez un cas où c'est sûr, et un cas où ce ne l'est pas.
\omcode{code/interviews-3/08-smart-contracts-and-crypto-engineering/solidity/Overflow.sol}{4}{16}{Dépassement de capacité et troncature.\protect\hfil\protect\break}\label{lst:iv3:cr:overflow}
\end{question}

\begin{question}[$\star\star$ \iqroles{QD, Dev} \iqfirm{crypto} \iqlevel{junior, confirmé}]\label{iq:iv3:smart-contracts-and-crypto-engineering:10}
Le \texttt{Vault} local, présenté dans le \cref{lst:iv3:cr:vault}, détient 10 ETH appartenant à Bob. La fonction \texttt{Attacker.attack()} est appelée avec 1 ETH. Combien détiennent le coffre et l'attaquant ensuite~? Expliquez la faille étape par étape. Donnez deux façons de la corriger, et dites laquelle vous choisiriez. Le retrait de Bob fonctionne-t-il encore après chaque correction~?
\omcode{code/interviews-3/08-smart-contracts-and-crypto-engineering/solidity/Vault.sol}{4}{33}{Un coffre et un appelant.\protect\hfil\protect\break}\label{lst:iv3:cr:vault}
\end{question}

\begin{question}[$\star\star$ \iqroles{QD, Dev} \iqfirm{crypto} \iqlevel{junior, confirmé}]\label{iq:iv3:smart-contracts-and-crypto-engineering:11}
\texttt{Payout}, présenté dans le \cref{lst:iv3:cr:pay}, paie un contrat qui refuse tous les paiements. Après \texttt{payUnchecked} avec 1 ETH, que montrent \texttt{paid}, le solde du destinataire et celui de \texttt{Payout}~? Que fait \texttt{payChecked}~? Quelles autres défaillances un \texttt{call} brut favorise-t-il~? Un contrat paie 100 destinataires dans une boucle~: pourquoi est-ce une mauvaise conception~?
\omcode{code/interviews-3/08-smart-contracts-and-crypto-engineering/solidity/Pay.sol}{4}{21}{Payer quelqu'un qui refuse.\protect\hfil\protect\break}\label{lst:iv3:cr:pay}
\end{question}

\begin{question}[$\star\star$ \iqroles{QD, Dev} \iqfirm{crypto} \iqlevel{junior}]\label{iq:iv3:smart-contracts-and-crypto-engineering:12}
Le propriétaire du \texttt{Wallet} local, présenté dans le \cref{lst:iv3:cr:origin}, qui détient 5 ETH, se laisse persuader d'appeler \texttt{Trap.claimPrize}. Que se passe-t-il lorsque \texttt{byOrigin} vaut vrai, et lorsqu'elle vaut faux~? Quelle est la différence entre \texttt{tx.origin} et \texttt{msg.sender}~? Pourquoi la vérification ne protège-t-elle pas le propriétaire~?
\omcode{code/interviews-3/08-smart-contracts-and-crypto-engineering/solidity/Origin.sol}{4}{28}{Qui appelle~?\protect\hfil\protect\break}\label{lst:iv3:cr:origin}
\end{question}

\begin{question}[$\star\star$ \iqroles{QD, Dev} \iqfirm{crypto} \iqlevel{junior, confirmé}]\label{iq:iv3:smart-contracts-and-crypto-engineering:13}
Dans un jeton standard, \texttt{approve(spender, m)} écrase l'allocation. Un propriétaire a accordé une allocation de 100 et envoie une transaction qui la ramène à 60. Le dépensier surveille les transactions en attente et peut faire exécuter les siennes en premier. Quel est le montant maximal qu'il peut prélever~? Montrez-le en Python. Comment la répareriez-vous, dans le contrat et côté client~?
\end{question}

\begin{question}[$\star\star$ \iqroles{QD, Dev, QT} \iqfirm{crypto} \iqlevel{junior, confirmé}]\label{iq:iv3:smart-contracts-and-crypto-engineering:14}
Un pool à produit constant avec des frais de 0.3~\% détient 1000 ETH et 2\,000\,000 dollars d'un stablecoin, chacun avec 18 décimales. Écrivez \texttt{amountOut} en Solidity avec des entiers. Combien de dollars permet d'acheter 1 ETH, et comment ce montant se répartit-il entre les frais et l'impact sur le prix~? Pourquoi la division arrondit-elle vers le bas, et au profit de qui~? Comment vérifiez-vous que le contrat concorde avec un modèle en Python~?
\end{question}

\begin{question}[$\star\star$ \iqroles{QD, Dev, Risque} \iqfirm{crypto} \iqlevel{junior, confirmé}]\label{iq:iv3:smart-contracts-and-crypto-engineering:15}
Un prêteur valorise son collatéral au prix spot d'un pool à produit constant peu liquide, dont les réserves sont de 1000 jetons et 1000 dollars, avec des frais de 0.3~\% prélevés sur l'entrée, et prête jusqu'à 50~\% de cette valeur. Un acteur qui détient 100 jetons, valant 100 dollars, peut emprunter 9000 dollars au sein d'une même transaction, sans frais, et verse la totalité dans le pool. En Python, modélisez la séquence et trouvez le profit. Quel est le profit si le prêteur utilise un prix juste à la place~? Qu'est-ce qui rend l'attaque possible, et quelles sont les défenses~?
\end{question}

\begin{question}[$\star\star$ \iqroles{QD, Dev} \iqfirm{crypto} \iqlevel{confirmé}]\label{iq:iv3:smart-contracts-and-crypto-engineering:16}
Au lieu de lire le prix spot d'un pool, un contrat peut conserver une somme cumulée du prix multiplié par le temps écoulé en secondes. En Python, écrivez cet accumulateur, ainsi que le prix moyen entre deux relevés. Si le prix passe de 1 à 100 pendant une seule des 100 secondes, quelle est la moyenne~? Que doit désormais faire un attaquant~? Quel compromis impose la fenêtre~? Où une moyenne arithmétique se comporte-t-elle mal ici~?
\end{question}

\begin{question}[$\star\star$ \iqroles{QD, Dev} \iqfirm{crypto} \iqlevel{junior, confirmé}]\label{iq:iv3:smart-contracts-and-crypto-engineering:17}
Un flux de prix local présenté dans le \cref{lst:iv3:cr:feed} a rapporté pour la dernière fois 200\,000 à l'instant $T$. \texttt{careful} est appelée avec une limite d'une heure, à $T + 60$ s, puis à $T + 7200$ s. Ensuite, le flux rapporte $-1$. Que renvoient \texttt{trusting} et \texttt{careful} dans chacun de ces cas~? Que devrait vérifier en plus un consommateur prudent~? Que doit faire le protocole lorsqu'il refuse un prix~?
\omcode{code/interviews-3/08-smart-contracts-and-crypto-engineering/solidity/Feed.sol}{4}{25}{Deux consommateurs d'un flux de prix.\protect\hfil\protect\break}\label{lst:iv3:cr:feed}
\end{question}

\begin{question}[$\star\star$ \iqroles{QD, QT} \iqfirm{crypto} \iqlevel{junior, confirmé}]\label{iq:iv3:smart-contracts-and-crypto-engineering:18}
Un pool à produit constant détient 1000 jetons et 2\,000\,000 dollars, avec des frais de 0.3~\% prélevés sur l'entrée. Un swap en attente dépense 20\,000 dollars, et il sera annulé s'il reçoit plus de 1~\% de moins que la cotation actuelle. En Python, modélisez un trader qui achète en premier et revend aussitôt après, et trouvez le plus grand premier achat qui laisse passer la victime, ainsi que son profit. Comment cela varie-t-il avec la limite~? Comment appelle-t-on cela, et que peut faire l'utilisateur pour s'en protéger~? Pourquoi ne pas fixer la limite à 0~?
\end{question}

\begin{question}[$\star\star$ \iqroles{QD, Dev} \iqfirm{crypto} \iqlevel{confirmé}]\label{iq:iv3:smart-contracts-and-crypto-engineering:19}
Selon la règle 16-et-4, un octet non nul des données d'une transaction coûte 16 gas, et un octet nul coûte 4. En Python, calculez le coût d'un \texttt{transfer(address, uint256)} standard encodé en ABI, portant sur 1 jeton à 18 décimales, vers une adresse sans octet nul, puis celui d'une version compactée~: un sélecteur d'un octet, l'adresse de 20 octets, et le montant dont on a retiré les zéros de tête. Pourquoi les rollups y tiennent-ils tant~? De quoi se composent les frais sur une couche 2, et qu'est-ce qui change dans le code que vous écrivez pour une telle couche~?
\end{question}

\begin{question}[$\star\star$ \iqroles{QD, Dev} \iqfirm{crypto} \iqlevel{junior, confirmé}]\label{iq:iv3:smart-contracts-and-crypto-engineering:20}
En Python, écrivez un processeur de dépôts pour une bourse. Il voit les blocs avec leurs parents, et il crédite un dépôt dès qu'il réunit $N$ confirmations sur la meilleure chaîne. Une réorganisation peut remplacer des blocs récents, et la même transaction peut réapparaître sur la nouvelle branche. Qu'est-ce qui garantit qu'un crédit n'a lieu qu'une seule fois~? Que doit-il se passer si une réorganisation va plus profond que $N$~? Comment le testez-vous~?
\end{question}

\begin{question}[$\star\star$ \iqroles{QD, Dev} \iqfirm{crypto} \iqlevel{junior, confirmé}]\label{iq:iv3:smart-contracts-and-crypto-engineering:21}
Une bourse paie les retraits depuis un seul compte chaud, en utilisant les nonces dans l'ordre. Une transaction est bloquée, car ses frais sont trop bas, et l'opérateur veut la renvoyer. En Python, écrivez un émetteur et une chaîne simulée. Pourquoi une nouvelle tentative doit-elle conserver son nonce, et ne pas en recevoir un nouveau~? Qu'est-ce qu'un nonce bloqué fait aux demandes qui le suivent~? Comment relevez-vous les frais~? Qu'est-ce qui doit être enregistré avant toute diffusion~?
\end{question}

\begin{question}[$\star\star\star$ \iqroles{QD, Dev} \iqfirm{crypto} \iqlevel{confirmé}]\label{iq:iv3:smart-contracts-and-crypto-engineering:22}
Dans le \cref{lst:iv3:cr:proxy}, quelqu'un appelle \texttt{setX(0xdead)} à travers le proxy. Où aboutit cette écriture~? Que renvoie ensuite \texttt{impl()}, et que fait alors un appel de \texttt{x()} à travers le proxy~? Expliquez \texttt{delegatecall}. Comment disposeriez-vous le stockage pour éviter cela~? Que doit respecter une implémentation mise à jour~?
\omcode{code/interviews-3/08-smart-contracts-and-crypto-engineering/solidity/Proxy.sol}{4}{22}{Un proxy et son implémentation.\protect\hfil\protect\break}\label{lst:iv3:cr:proxy}
\end{question}

\begin{question}[$\star\star\star$ \iqroles{QD, Dev} \iqfirm{crypto} \iqlevel{confirmé}]\label{iq:iv3:smart-contracts-and-crypto-engineering:23}
Dans le \cref{lst:iv3:cr:readonly}, Bob a rejoint le pool avec 100 ETH, et \texttt{Reader.run()} est appelée avec 100 ETH de plus, sur une chaîne locale. Quel prix \texttt{Reader} voit-il dans \texttt{receive}, et quel est-il juste avant et juste après~? Aucun état du pool n'est modifié pendant le rappel, alors pourquoi s'agit-il d'une forme de réentrance~? Pourquoi un verrou sur \texttt{leave} seul ne suffit-il pas~? Que perdrait un prêteur qui lirait ce prix~? Corrigez-le.
\omcode{code/interviews-3/08-smart-contracts-and-crypto-engineering/solidity/ReadOnly.sol}{4}{36}{Un prix lu au milieu d'un appel.\protect\hfil\protect\break}\label{lst:iv3:cr:readonly}
\end{question}

\begin{question}[$\star\star\star$ \iqroles{QD, Dev} \iqfirm{crypto} \iqlevel{confirmé}]\label{iq:iv3:smart-contracts-and-crypto-engineering:24}
Le \Cref{lst:iv3:cr:relay} permet à n'importe qui de soumettre le paiement signé d'un utilisateur et d'en payer le gas. Le signataire a approvisionné 10 ETH dans chaque déploiement, et une signature paie 1 ETH à Mallory. Si Mallory la soumet deux fois, puis une fois à un autre déploiement du même contrat, que se passe-t-il~? Qu'est-ce qui manque au message signé, et pourquoi chaque champ manquant compte-t-il~? Écrivez la correction. Pourquoi un utilisateur peut-il payer sans détenir de gas de cette façon~?
\omcode{code/interviews-3/08-smart-contracts-and-crypto-engineering/solidity/Relay.sol}{4}{18}{Un paiement relayé.\protect\hfil\protect\break}\label{lst:iv3:cr:relay}
\end{question}

\begin{question}[$\star\star\star$ \iqroles{QD, Dev, Risque} \iqfirm{crypto} \iqlevel{confirmé}]\label{iq:iv3:smart-contracts-and-crypto-engineering:25}
Une clé qui contrôle des fonds doit survivre à la perte d'un détenteur, et résister au vol de l'un d'entre eux. En Python, écrivez le partage de secret de Shamir, $t$ parmi $n$, sur un corps premier. Pourquoi $t - 1$ parts ne révèlent-elles rien~? Si chacun des 3 détenteurs de clé est compromis indépendamment avec une probabilité de 1~\%, quelle est la probabilité d'un vol sous un schéma 2 sur 3, et quelle est-elle avec une seule clé~? En quoi cela diffère-t-il d'un contrat multisignature, et d'une signature à seuil~?
\end{question}

\begin{question}[$\star\star\star$ \iqroles{QD, Dev, Risque} \iqfirm{crypto} \iqlevel{confirmé}]\label{iq:iv3:smart-contracts-and-crypto-engineering:26}
Une bourse publie la racine d'un arbre de Merkle des soldes de comptes, afin que chaque utilisateur puisse vérifier le sien. En Python, écrivez l'arbre, une preuve d'inclusion et sa vérification. Pourquoi les feuilles et les nœuds internes reçoivent-ils des préfixes différents~? Quelle est, au plus, la longueur d'une preuve pour 23 comptes~? Que démontre-t-elle, et que ne démontre-t-elle pas~? Que faut-il encore montrer pour soutenir une affirmation de solvabilité~?
\end{question}

\begin{question}[$\star\star\star$ \iqroles{QD, QT, Risque} \iqfirm{crypto} \iqlevel{confirmé}]\label{iq:iv3:smart-contracts-and-crypto-engineering:27}
Une position longue perpétuelle de 1 jeton à 2000, avec une marge isolée de 200 et un taux de maintenance de 0.5~\%, est clôturée lorsque les fonds propres tombent au niveau du taux de maintenance multiplié par le notionnel évalué au prix de marquage. En Python, trouvez le prix de liquidation. Pourquoi utiliser un prix de marquage plutôt que la dernière transaction~? Le prix saute, et la clôture s'exécute à 1790~: qui paie~? Que se passe-t-il si le fonds d'assurance est épuisé~?
\end{question}

\begin{question}[$\star\star\star$ \iqroles{QD, QT} \iqfirm{crypto} \iqlevel{confirmé}]\label{iq:iv3:smart-contracts-and-crypto-engineering:28}
Un pool à produit constant détient 1000 jetons et 2\,000\,000 dollars, avec des frais de 0.3~\% prélevés sur l'entrée, et une plateforme de négociation centralisée affiche une cotation de 2020 pour le jeton. En négligeant les coûts de cette plateforme, dérivez la taille optimale d'achat dans le pool et de revente sur cette plateforme, puis trouvez le profit. En dessous de quel prix externe n'y a-t-il aucune transaction~? Écrivez-la en Python. Quels risques la formule omet-elle, à l'exécution et dans l'inventaire~?
\end{question}

\begin{question}[$\star\star\star\star$ \iqroles{QD, Dev, LLE} \iqfirm{crypto} \iqlevel{confirmé}]\label{iq:iv3:smart-contracts-and-crypto-engineering:29}
Concevez une bourse crypto centralisée avec dépôts, retraits, ordres à cours limité et carnet d'ordres en direct. Comment sont tenus les soldes, et comment sont-ils bloqués pour les ordres ouverts~? En Python, écrivez le noyau des soldes. Où le moteur d'appariement rencontre-t-il le registre~? Comment fonctionnent les dépôts, les retraits et la garde des actifs, et quelle part des fonds est placée dans un portefeuille chaud~? Qu'est-ce qui ne doit jamais être possible, et comment vérifieriez-vous que cela ne s'est jamais produit~?
\end{question}

\begin{question}[$\star\star\star\star$ \iqroles{QD, QT, LLE} \iqfirm{crypto} \iqlevel{confirmé}]\label{iq:iv3:smart-contracts-and-crypto-engineering:30}
Concevez un système qui fait du market making sur un pool on-chain et couvre ses positions sur des plateformes de négociation centralisées. Que lit-il sur la chaîne, et à quelle étape d'un bloc~? Comment les réorganisations, les transactions échouées et le gas modifient-ils la conception~? Comment gérez-vous le délai entre une décision et son inclusion~? Dans quel état se trouve une transaction que vous avez envoyée, et comment votre position est-elle suivie~? Que surveilleriez-vous~? Que doit savoir faire un environnement de test~?
\end{question}
